\documentclass[10pt,a4paper]{amsart}
\usepackage[margin=19mm]{geometry}
\usepackage{amsmath,amssymb,mathtools}
\usepackage[scr=rsfs,cal=euler]{mathalfa}
\usepackage{xfrac}
\usepackage[unicode,hidelinks,bookmarks=false]{hyperref}
\newcommand{\ie}{i.e.\ }
\newcommand{\cf}{cf.\ }
\newcommand{\resp}{respectively\ }
\newcommand{\Subsection}{\subsection}
\providecommand{\nobreakdash}{\nobreak\hbox{-}}
\setlength{\emergencystretch}{2em}
\multlinegap0pt
\usepackage{xcolor}
\usepackage[shortlabels]{enumitem}
\usepackage{amssymb}
\usepackage{mathtools}
\newtheorem{prop}{Proposition}
\usepackage{cleveref}
\crefname{prop}{Proposition}{Propositions}
\newcommand{\qop}{\triangleleft}
\title{Detecting non-admissibility of quandles via colorings}
\author{Katsunori Arai, Ryoya Kai}
\date{Source: arXiv:2603.29320v1 (CC BY 4.0)}
\begin{document}
\maketitle

\noindent\emph{Source and licence.} Detecting non-admissibility of quandles via colorings. Version arXiv:2603.29320v1 (CC BY 4.0), distributed by arXiv under the Creative Commons Attribution 4.0 International licence, http://creativecommons.org/licenses/by/4.0/, as recorded in the arXiv licence metadata for this exact version and shown on the abstract page https://arxiv.org/abs/2603.29320v1. Full licence notice: \texttt{licenses/arxiv-2603.29320-CC-BY-4.0.txt}.\par\smallskip

\noindent\textit{Editorial scope.} Counted unit: the article's prop Prop:HopflinkNAdm (source0.tex lines 269--272). The article's complete original proof of it is delivered with it (source0.tex lines 273--298, one \texttt{proof} environment of 26 lines). No mathematics is written, repaired or re-derived: the statement and the proof are the article's own lines, in the article's own notation, and no retained line is substituted or re-flowed.  No further source-local context is needed: every cross-reference of every retained line resolves inside the two delivered spans.  One declared adaptation: exactly 1 ancillary strings inside the retained spans are omitted (image-inclusion commands and, where the source repeats a label, the repeated label definitions) and no other line differs from the source; the figure environments, with their placement, captions and labels, are kept, and the package records the omitted strings in fidelity receipts bound to the affected bodies.  No retained line contains a citation, so the wrapper carries no bibliography and no bibliography entry.  The retained lines contain the article's own diagram code, which the wrapper typesets with the standard tikz packages; no image file is delivered and the package declares no asset dependency.  Editorial formatting only, in addition: an AMS \texttt{amsart} preamble that restates the article's own declarations and macros used by the retained lines, namely the environments \texttt{prop} on separate counters, together with the standard packages those lines need.  The retained environments print in this document as Proposition 1 in order of appearance, whereas the article prints the same items with the numbering of its own full document.  The complete native source of the article as distributed in the arXiv e-print for this exact version is bundled unchanged as \texttt{sources/197507/source0.tex}.
\begin{prop}
  A quandle $X$ is not $2_{1}^{2}$-admissible if and only if there exist $x, y \in X$ such that $x \qop y = x$ and $y \qop x \neq y$.
  \label{Prop:HopflinkNAdm}
\end{prop}

\begin{proof}
  Let $D$ be the diagram of the $(1,1)$-tangle obtained from the Hopf link $2_{1}^{2}$ as shown in \cref{Fig:hopftanglediagram}.
  Then, any $X$-coloring of $D$ is determined as follows:
  first,
  we assign two elements $x, y \in X$ to the arcs as in \cref{Fig:hopftanglediagram}.
  Then, the element assigned to the arc $a_{t}$ is determined by the relation at the crossing $c_{1}$.
  Finally, if the elements $x$ and $y$ satisfy the relation at $c_{2}$, which is given by
  \begin{equation*}
    x \qop (y \qop x) = x \Leftrightarrow x \qop y = x,
  \end{equation*}
  then the map $C: \mathrm{Arc}(D) \to X$ is an $X$-coloring.
  In addition,
  it satisfies
  \begin{equation*}
    C(a_{s}) \neq C(a_{t}) \Leftrightarrow y \qop x \neq y.
  \end{equation*}
  We note that the above condition does not depend on the choice of an orientation of $2_{1}^{2}$ and a base point.
  Therefore,
  the quandle $X$ satisfies the condition of this proposition if and only if $X$ is not $2_{1}^{2}$-admissible.  
  \begin{figure}[h]
    \centering
    
    \caption{An $X$-colored $(1,1)$-tangle diagram $D$ obtained from a diagram of the Hopf link $2_{1}^{2}$}
    \label{Fig:hopftanglediagram}
  \end{figure}
\end{proof}

\end{document}
